\section{年度 Aggregate Risk 模拟}

\subsection{模拟框架}

在频率分布、严重度分布与状态转移矩阵都确定之后，项目进入年度风险模拟阶段。其目标不是预测某个具体月，而是给出 \textbf{未来一年总 delay-equivalent minutes} 的分布。

设一年有 12 个月。对第 $m$ 个月，先生成状态 $S_m \in \{\text{normal}, \text{disrupted}\}$；再随机抽取一个月度样本中的航司数量 $A_m$；对每家航司 $a=1,\dots,A_m$，按选中的频率模型生成 internal/system/external/cancellation/diversion 事件数，再按选中的严重度模型生成每类事件的平均分钟损失。

于是月度总冲击可以写为
\begin{equation}
L_m
=
\sum_{a=1}^{A_m}
\left(
\sum_{k=1}^{N_{m,a}^{(\text{int})}} X_{m,a,k}^{(\text{int})}
\;+\;
\sum_{k=1}^{N_{m,a}^{(\text{sys})}} X_{m,a,k}^{(\text{sys})}
\;+\;
\sum_{k=1}^{N_{m,a}^{(\text{ext})}} X_{m,a,k}^{(\text{ext})}
\;+\;
180\,C_{m,a}
\;+\;
240\,V_{m,a}
\right).
\end{equation}

因此年度总冲击为
\begin{equation}
L_{\text{annual}} = \sum_{m=1}^{12} L_m.
\end{equation}

项目设置了 500 次 Monte Carlo 模拟，并固定随机种子为 7337，以保证复现性。

\subsection{状态驱动的模拟步骤}

脚本中的模拟顺序可概括为：
\begin{enumerate}[leftmargin=2.2em]
  \item 使用 Markov 链或状态占比生成 12 个月的状态序列；
  \item 根据状态选择对应的情景乘子；
  \item 对每个月抽取月度航司数；
  \item 对每家航司，按照选中的频率模型生成 internal/system/external/cancel/divert 五类事件数；
  \item 对 internal/system/external 三类事件，按照选中的严重度模型生成平均事件损失并乘以事件数；
  \item 对取消与备降直接施加 180 与 240 分钟 penalty；
  \item 汇总得到年度总冲击，并在重复模拟后计算均值、VaR 与 TVaR。
\end{enumerate}

\subsection{情景乘子设计}

项目不是只做单一 baseline，而是引入了 4 类运营情景：
\begin{itemize}[leftmargin=2.2em]
  \item \texttt{base}
  \item \texttt{disruption\_stress}
  \item \texttt{weather\_shock}
  \item \texttt{holiday\_peak}
\end{itemize}

这些情景通过对频率与严重度施加乘子实现。若某一风险块基准频率均值为 $\mu$，情景乘子为 $m$，则模拟中使用的调整后均值为
\begin{equation}
\mu^{\star} = m\mu.
\end{equation}

对严重度也使用同样思路：
\begin{equation}
X^{\star} = mX.
\end{equation}

下表给出脚本中的乘子设置。为便于阅读，仅保留仓库中实际使用的数值。

\begin{landscape}
\begin{table}[H]
\centering
\caption{正常状态下的情景乘子}
\resizebox{\linewidth}{!}{
\begin{tabular}{lcccccccc}
\toprule
Scenario & internal & system & external & cancel & divert & sev\_internal & sev\_system & sev\_external \\
\midrule
base & 1.00 & 1.00 & 1.00 & 1.00 & 1.00 & 1.00 & 1.00 & 1.00 \\
disruption\_stress & 1.05 & 1.10 & 1.10 & 1.05 & 1.05 & 1.03 & 1.05 & 1.05 \\
weather\_shock & 1.00 & 1.05 & 1.25 & 1.10 & 1.08 & 1.00 & 1.03 & 1.12 \\
holiday\_peak & 1.12 & 1.12 & 1.08 & 1.08 & 1.05 & 1.05 & 1.05 & 1.03 \\
\bottomrule
\end{tabular}}
\end{table}

\begin{table}[H]
\centering
\caption{扰动状态下的情景乘子}
\resizebox{\linewidth}{!}{
\begin{tabular}{lcccccccc}
\toprule
Scenario & internal & system & external & cancel & divert & sev\_internal & sev\_system & sev\_external \\
\midrule
base & 1.10 & 1.20 & 1.25 & 1.15 & 1.10 & 1.05 & 1.08 & 1.10 \\
disruption\_stress & 1.25 & 1.35 & 1.45 & 1.30 & 1.20 & 1.10 & 1.15 & 1.20 \\
weather\_shock & 1.05 & 1.15 & 1.65 & 1.40 & 1.30 & 1.02 & 1.08 & 1.25 \\
holiday\_peak & 1.22 & 1.28 & 1.18 & 1.22 & 1.15 & 1.10 & 1.12 & 1.08 \\
\bottomrule
\end{tabular}}
\end{table}
\end{landscape}

\subsection{年度风险指标}

项目输出以下风险指标：
\begin{itemize}[leftmargin=2.2em]
  \item \textbf{Expected impact}：模拟样本均值；
  \item \textbf{VaR 95\%}：95\% 分位数；
  \item \textbf{TVaR 95\%}：超过 VaR 95\% 后的条件均值；
  \item \textbf{VaR 99\%}：99\% 分位数；
  \item \textbf{TVaR 99\%}：超过 VaR 99\% 后的条件均值。
\end{itemize}

若 $L$ 是年度总冲击，则
\begin{equation}
\mathrm{VaR}_{\alpha}(L) = \inf \{x : \mathbb{P}(L \le x) \ge \alpha\},
\end{equation}
而
\begin{equation}
\mathrm{TVaR}_{\alpha}(L) = \mathbb{E}[L \mid L \ge \mathrm{VaR}_{\alpha}(L)].
\end{equation}

\subsection{模拟结果}

\begin{table}[H]
\centering
\caption{年度 aggregate risk 情景结果}
\resizebox{\textwidth}{!}{
\begin{tabular}{lrrrrr}
\toprule
情景 & Expected impact & VaR 95\% & TVaR 95\% & VaR 99\% & TVaR 99\% \\
\midrule
base & 2{,}367{,}072.82 & 2{,}790{,}834.28 & 2{,}911{,}906.09 & 3{,}000{,}309.69 & 3{,}031{,}286.85 \\
disruption\_stress & 2{,}611{,}594.29 & 3{,}121{,}850.75 & 3{,}323{,}801.55 & 3{,}448{,}514.00 & 3{,}536{,}246.17 \\
weather\_shock & 2{,}455{,}662.60 & 2{,}884{,}871.28 & 3{,}016{,}874.51 & 3{,}066{,}805.62 & 3{,}234{,}528.83 \\
holiday\_peak & 2{,}672{,}396.40 & 3{,}182{,}046.44 & 3{,}277{,}558.29 & 3{,}323{,}075.10 & 3{,}417{,}469.59 \\
\bottomrule
\end{tabular}}
\end{table}

\begin{figure}[H]
  \centering
  \includegraphics[width=0.84\textwidth]{chart_6_aggregate_risk_scenarios.png}
  \caption{不同情景下的年度 aggregate risk 分布}
  \label{fig:scenario-kde}
\end{figure}

对这些结果可以做如下解读：
\begin{itemize}[leftmargin=2.2em]
  \item \texttt{base} 情景下，年度期望冲击约为 236.7 万分钟；
  \item \texttt{holiday\_peak} 的期望冲击最高，说明高峰期流量上升会显著推高年度总损失；
  \item \texttt{disruption\_stress} 的尾部指标最重，说明该情景更能放大极端年份的风险；
  \item \texttt{weather\_shock} 虽然均值不如前两者高，但在尾部仍然保持较高 TVaR，说明外部冲击主要体现为尾部风险扩张。
\end{itemize}

\subsection{这一阶段的统计含义}

从统计上看，项目的 aggregate risk 模块完成了如下整合：
\begin{enumerate}[leftmargin=2.2em]
  \item 用 Negative Binomial 描述事件数量的过度离散；
  \item 用 Lognormal / Weibull 描述事件严重度的右偏分布；
  \item 用 Markov 链把 ``月份环境'' 的持续性显式纳入模拟；
  \item 用 scenario multipliers 把管理者关心的运营情景映射为可计算的分布变化。
\end{enumerate}

因此，最终的年度风险不是单一分布假设的产物，而是 \textbf{频率、严重度、状态依赖和情景冲击的复合输出}。
